\documentclass{article}

\usepackage{graphicx}

\title{Stenosis articles}
\author{Dmitrii Sakharov}
\date{March 2026}

\begin{document}

\maketitle

\section{Introduction}
\cite{aipoweredrealtime}. Task of segmentation and detection united, keyframe selection (keyframe selection is done by edge detection followed up with an image processing and calculation of vessel area. First, there is a segformer that segments the whole tree and find different parts (rca, pda, etc.). Then rcnn finds boxes with leison, then those boxes + surroundings are binary segmented to get a nice mask and then skeletonize and geometrically calculate diameter. 

\cite{fvcm}. Here they choose the best frames with CNN+LSTM, the best frames they calssify (just put the whole frame into the classifier). Then they use different projections from patietns to max-pool so to check whether any of the frames are marked "stenosis". Then they use Anchor-based FPN (Feature Pyramid Network) for detection. They had a problem of detection many small leisons in small branches. 

\cite{molenar}. Here they segment eveything into lad+lcx+rca (multicalss segmentation) with resnet101 and u-net. Then detected edges of the vessel, calculated qca. Limitations: single frames + only main arteries + obly Phillips. 


\cite{popov2022reviewmodernapproachescoronary}. Overview. Research gaps: no stnardatized normal metrics/data/benchmarks, low generalizability, not enough external validation. Ignoring of small branches, not enough video-based models. No model confidence metric for stenosis.  

\cite{ZHAO2021104667}. Here they binary segment the whole tree (FP-U-Net++). Then they calculate central line, qca. Limiatations: single frame. 

\cite{cardiovascular}. Just cardiovasular review, but points out a bug problem with domain shift (differnt x ray machines).

\cite{attentionstenosis}. Used several objcet detection models to detect stenosis (binary detection). Used it only on right arteries. The best model was EfficientDet (D3) due to the dual attention of BiFPN! Still had many FPs, used only single frames, so segmentation.

\cite{MOON2021105819}. Here they use a weakly-supervised learning approach. Keyframe selection is done automatically using a top-hat transform and Frangi filter to find the frame with the highest number of vessel pixels. Then they classify the whole frame (normal vs. stenosis $\ge50\%$) using GoogleNet Inception-v3 combined with CBAM (spatial and channel attention module). For localization, they use Grad-CAM to generate a heatmap of the stenosis, avoiding the need for bounding box or mask annotations during training. Limitations: only right coronary artery (RCA), rough localization (heatmap instead of precise QCA), and critical dependence on keyframe selection (error rate spikes if the algorithm selects the wrong frame).

\cite{STQD-Det}. Here they tackle the problem at the video level (video object detection) by inputting a sequence of 9 consecutive frames, bypassing explicit keyframe selection. They propose a Quantum Diffusion Model (using discrete Poisson noise instead of Gaussian to generate bounding boxes) with a ResNet50 + FPN backbone. The core innovation is the Spatio-Temporal Feature Sharing (STFS) module. It uses the Hungarian algorithm to match detected boxes across frames and projects features (RoI) from "confident" (positive) frames onto problematic/error frames, fixing missed or false detections caused by vessel motion or poor contrast. Limitations: only bounding box detection (localization), no quantitative evaluation of stenosis (QCA), and no severity grading/classification.

\cite{menezes2023}. A multicentric validation study of a deep learning model for multi-class segmentation (background, coronary tree, and catheter). It utilizes the EfficientUNet++ architecture with an EfficientNet-B5 encoder. The model was tested on a dataset from 4 different medical centers (123 measurements) and compared against commercial QCA software (CAAS Workstation) as the ground truth. The model achieved a Dice Score of 94.8, an IoU of 90.1, and a clinician-based Global Segmentation Score (GSS) of 92/100. Statistical analysis showed no significant differences in the calculation of lesion diameter or percentage diameter stenosis (DS) between the AI model and manual QCA. Limitations: the study is limited to single-frame analysis, features a relatively small sample size, focuses only on diseased segments rather than the full coronary tree, and excludes images with artifacts such as pacemakers or metallic clips. SEGEMNTATION NOT DETECTION!

\cite{HAN2023106546}. This paper proposes a PS-STT (Proposal-Shifted Spatial-Temporal Transformer) for stenosis detection in X-ray sequences. It bypasses single-frame limitations by processing 9-frame sequences with a ResNet-50 + FPN backbone. The core innovation is the proposal-shifted mechanism, which aligns vessel positions across time to handle cardiac motion. It outperforms state-of-the-art video object detection models with an mAP of 78.47 and an F1-score of 79.11. Limitations: it provides only bounding box localization (no precise QCA or masks) and focuses on detection rather than severity grading.

\cite{huang2026vasomim}. Focuses on self-supervised pre-training (SSL) using the VasoMIM (Vascular anatomy-aware Masked Image Modeling) framework to build a foundation model for X-ray angiograms. It explicitly incorporates vascular anatomy into the learning process through an anatomy-guided masking strategy (prioritizing the masking of informative vessel patches over background) and an anatomical consistency loss to preserve vascular topology during reconstruction. The model was pre-trained on XA-170K, currently the largest dataset in the domain (171,478 images). It achieves SOTA across four downstream tasks: vessel segmentation, vessel segment segmentation, stenosis segmentation, and stenosis detection. It demonstrates extreme label efficiency, where fine-tuning on only 10percent of labels outperforms fully-supervised models trained on 100percent of data. Limitations: limited to single-frame analysis, marginal returns when scaling model capacity (ViT-L/H) or data volume above 40K frames, and potential architectural bottlenecks in the downstream segmentors. BOTH SEGMENTATION AND DETECTION!

\cite{JUNGIEWICZ2023120234}. This paper compares different variants of the Inception-v3 Network (CNN) and the Vision Transformer (ViT) for the binary classification of stenosis in coronary angiography. The study utilizes small image patches (32x32 pixels) rather than full frames. They introduce a new dataset of ~2400 samples from the University Hospital in Krakow and address artifacts found in existing public datasets, such as image borders. A key experiment involves analyzing model performance across varying percentages of arteries in negative samples (0percent to 100percent) to ensure the model detects stenosis rather than just the presence of a vessel. To improve ViT performance on a relatively small medical dataset, they employ Sharpness-Aware Minimization (SAM). Grad-CAM is used to provide visual explanations of model focus. Results show that while CNNs generally perform better, the SAM-HVITFT (Half-ViT with fine-tuning and SAM) is the most stable model and achieves comparable performance to CNNs, even outperforming them in mean F1 score in some scenarios. Limitations: The study is limited to single-frame (patch-based) analysis and does not account for temporal information in video sequences. Additionally, the Transformer models were constrained by a small patch grid ($4\times4$) due to the small input size, and the overall dataset remains relatively small for large-scale Transformer training. CLASSIFICATION NOT DETECTION/SEGMENTATION!

\cite{zohranyan2024drsamendtoendframeworkvascular}. This paper presents Dr-SAM, an end-to-end framework for vascular segmentation, diameter estimation, and anomaly detection in peripheral angiography. The method leverages the Segment Anything Model (SAM) by introducing a specialized positive point selection mechanism that allows for precise vessel extraction without the need for fine-tuning or additional training. After segmentation, the framework extracts a topological skeleton, prunes noise, and uses distance transforms to estimate vessel diameters and identify anomalies (stenosis and aneurysms) as extremum points. The study introduces a new benchmark dataset of 500 images of pelvic-iliac arteries with expert-validated binary masks and anomaly labels. Limitations: The current pipeline requires a user-provided bounding box for initial point selection. It is specifically optimized for large peripheral vessels (aorta and iliac arteries) rather than the entire coronary tree. SEGMENTATION, DIAMETER ESTIMATION, AND ANOMALY DETECTION!

\cite{lin2025deepselfknowledgedistillationhierarchical}. This paper introduces a Deep self-knowledge distillation (Deep self-KD) method for coronary artery segmentation in X-ray angiography. Instead of relying on a pre-trained static teacher, the student model learns from its own past iterations (the teacher model) using a hierarchical learning strategy on an encoder-decoder architecture (primarily U-Net3+). The core innovation is the combination of two losses: a deep distribution loss that applies a "loosely constrained" KL divergence on probabilistic distribution vectors derived from multi-level side outputs, and a pixel-wise self-KD loss that uses a linearly scheduled soft target (teacher output + ground truth) to regularize final predictions. Evaluated on the XCAD and DCA1 datasets, it achieved state-of-the-art Dice scores (80.88percent and 81.06percent, respectively) without adding any computational cost during inference. Limitations: single-frame analysis only, moderate dataset sizes (small-sample robustness is unproven), lack of investigation into how the model handles noisy or inaccurately annotated labels, and unverified generalizability to other imaging modalities. SEGMENTATION NOT DETECTION!


\cite{PU2023106493} Semi-supervised segmentation of coronary DSA using mixed networks and multi-strategies. 

\cite{JUNGIEWICZ2023120234}. This paper compares Inception-v3 (CNN) and Vision Transformer (ViT) variants for the binary classification of stenosis on small 32x32 coronary angiography patches. Using a custom dataset of ~2400 samples, they evaluate performance across varying percentages of arteries in negative samples to ensure the model targets stenosis rather than just vessel presence. To overcome the ViT's typical reliance on massive datasets, the study focuses on applying Sharpness-Aware Minimization (SAM). By simultaneously optimizing both the loss value and its sharpness, SAM guides the ViT to a flatter minimum, significantly boosting its stability and generalization on limited medical data. While CNNs generally perform better overall, the SAM-HVITFT (Half-ViT with fine-tuning and SAM) proves to be the most stable model against data variations and ultimately outperforms the CNNs in the mean F1 score. Limitations: restricted to static patch-based analysis without temporal video context and constrained by a small 4x4 patch grid due to the input size. CLASSIFICATION NOT DETECTION/SEGMENTATION!

\cite{PANG2021101900}. Here they propose Stenosis-DetNet, tackling stenosis detection using a sequence of 9 consecutive frames[cite: 74, 88]. The model extracts candidate boxes using a ResNet-50 + FPN backbone[cite: 101, 102]. The core innovations are two modules: Sequence Feature Fusion (SFF) and Sequence Consistency Alignment (SCA)[cite: 98]. SFF fuses candidate box features across frames based on similarity, enhancing stenosis characteristics against motion and contrast artifacts[cite: 104, 106]. SCA optimizes the results by clustering detected boxes across frames using IoU and structural similarity (SSIM); it removes false positives that do not appear consistently across a set number of frames and uses linear interpolation to recover missed bounding boxes on intermediate frames[cite: 211, 231, 234, 238]. Evaluated on 166 clinical sequences, it achieved an F1-score of 88.1\%[cite: 85, 511]. Limitations: it provides only bounding box localization without precise vessel segmentation or quantitative diameter calculations, and it currently relies on single-angle X-ray radiography[cite: 75, 571]. DETECTION NOT SEGMENTATION!



\bibliographystyle{plain} 
\bibliography{stenosis} 

\end{document}